%=======================================================================
%  Phase II Progress Report — HCM Variant Pathogenicity Classification
%  RVCE Experiential Learning Project | Academic Year 2025-26
%=======================================================================
\documentclass[12pt,a4paper]{article}

%---------- Packages ----------
\usepackage[T1]{fontenc}
\usepackage[utf8]{inputenc}
\usepackage[margin=1in]{geometry}
\usepackage{booktabs}
\usepackage{graphicx}
\usepackage{amsmath}
\usepackage{amssymb}
\usepackage{hyperref}
\usepackage{array}
\usepackage{longtable}
\usepackage{multirow}
\usepackage{xcolor}
\usepackage{caption}
\usepackage{subcaption}
\usepackage{setspace}
\usepackage{enumitem}
\usepackage{float}
\usepackage{lmodern}
\usepackage{parskip}

\hypersetup{
    colorlinks=true,
    linkcolor=blue!60!black,
    citecolor=blue!60!black,
    urlcolor=blue!60!black,
}

\setstretch{1.25}

%---------- Title block ----------
\title{%
  \LARGE\bfseries In Silico Prioritization of HCM Variants:\\
  A Comparative Study of Tabular and\\
  Sequence-Based Hybrid Modeling\\[1.2em]
  \large Phase II Progress Report\\
  \normalsize RVCE Experiential Learning Programme
}
\author{%
  Abhilash S.\\[0.3em]
  \small Department of Biotechnology\\
  \small RV College of Engineering, Bengaluru
}
\date{April 2026}

%=======================================================================
\begin{document}
\maketitle
\thispagestyle{empty}
\newpage

%---------- Table of Contents ----------
\tableofcontents
\newpage

%=======================================================================
\begin{abstract}
Hypertrophic Cardiomyopathy (HCM) is the most prevalent inherited cardiac
disorder, affecting approximately 1 in 500 individuals worldwide. Despite
rapid advances in clinical genomic sequencing, a substantial proportion of
identified missense variants remain classified as \emph{Variants of Uncertain
Significance} (VUS), creating a diagnostic gap that impedes genetic
counselling and family-cascade screening. This report documents the design,
execution, and evaluation of a Phase~II comparative machine-learning
investigation to address this gap.

A dataset of \textbf{1,954 clinically annotated missense variants} spanning
nine sarcomeric genes was assembled from ClinVar, ClinGen, gnomAD, ExAC,
TOPMed, dbSNP, and UniProt. Two primary modelling streams were constructed
and rigorously compared: (1)~a \emph{Tabular} stream powered by a tuned
XGBoost ensemble operating on 29 hand-crafted biochemical and structural
features, and (2)~a \emph{Sequence} stream powered by a one-dimensional
Convolutional Neural Network (1D-CNN) operating on an 11-residue
sliding-window encoding of the local amino-acid environment. A
\emph{Late-Fusion Hybrid} stacking both streams via a Logistic Regression
meta-learner constitutes the third experimental branch.

All models were evaluated under 5-Fold Stratified Cross-Validation to
respect the 73.7\,\% pathogenic class imbalance. The Hybrid model achieved
the best overall performance with an ROC-AUC of \textbf{0.8222}, a
Matthews Correlation Coefficient (MCC) of \textbf{$\approx$0.45}, and a
Balanced Accuracy of \textbf{0.753}, outperforming both individual streams
across all balanced metrics. These findings validate the complementarity
hypothesis and establish a strong foundation for Phase~III external
validation using the SHaRe registry and protein language model embeddings.
\end{abstract}
\newpage

%=======================================================================
\section{Introduction}
\label{sec:intro}

\subsection{The VUS Crisis in Inherited Cardiomyopathy}
\label{subsec:vus}

Hypertrophic Cardiomyopathy arises predominantly from heterozygous missense
mutations in genes encoding sarcomeric proteins. The clinical severity of
HCM is highly variable: some mutation carriers remain asymptomatic for
decades, whereas others experience sudden cardiac death in early adulthood.
Accurate identification of disease-causing variants is therefore not merely
an academic exercise but a matter of direct clinical consequence.

Modern next-generation sequencing panels routinely interrogate the nine
principal HCM-associated genes—\textit{MYH7}, \textit{MYBPC3},
\textit{TNNT2}, \textit{TNNI3}, \textit{TPM1}, \textit{ACTC1},
\textit{MYL2}, \textit{MYL3}, and \textit{TNNC1}—producing thousands of
variant calls per proband. The critical limitation is classification
confidence: a large fraction of missense variants discovered in clinical
panels are assigned VUS status by human curators, particularly for
less-studied genes. The 2015 ACMG/AMP guidelines provide a framework for
variant interpretation, but their application remains labour-intensive,
requires specialist structural knowledge, and is inconsistent across
laboratories.

The downstream consequences are substantial. VUS reports leave clinicians
unable to issue actionable recommendations to patients and families. They
inflate the psychological burden on carriers who receive ambiguous results.
They delay prophylactic interventions—implantable cardioverter-defibrillators,
myomectomy, mavacamten therapy—that could prevent adverse events.
Computational tools that can convert VUS into high-confidence predictions
therefore address a real and urgent unmet need in cardiology genomics.

\subsection{The Hybrid Modelling Rationale}
\label{subsec:hybrid}

Prior computational approaches to variant pathogenicity prediction fall into
two broad camps. \emph{Sequence-agnostic} methods—SIFT, PolyPhen-2, CADD,
and their successors—use conservation and physicochemical features but do
not model the local sequence neighbourhood explicitly. \emph{Deep-learning}
methods—VARITY, ClinPred, EVE—leverage sequence information but often
require large protein-family multiple sequence alignments and are not
designed to exploit gene-level structural annotation databases such as
UniProt.

Our Phase~II hypothesis is that these two evidence streams are
\emph{complementary rather than redundant}. Structured biochemical and
structural features encode expert knowledge accumulated over decades of
protein biochemistry and clinical genetics. Sequence-derived features
encode neighbourhood-level physicochemical context that is invisible to
tabular methods. A hybrid architecture that fuses both information types
should, in principle, outperform either stream in isolation—particularly on
the minority Benign class, where tabular models exhibit systematic recall
failure under class imbalance.

%=======================================================================
\section{Technical Methodology}
\label{sec:methods}

\subsection{Dataset Curation and Class Distribution}
\label{subsec:data}

The curated dataset comprises \textbf{1,954 missense variants} drawn
exclusively from nine HCM-associated sarcomeric genes. All variants carry
binary ground-truth labels: Pathogenic ($y=1$) or Benign ($y=0$), derived
from ClinVar significance annotations. Variants of Uncertain Significance
(VUS) were withheld from the supervised training set to preserve label
integrity. Table~\ref{tab:gene_dist} documents the per-gene breakdown.

\begin{table}[H]
\centering
\caption{Per-gene variant count and proportion in the curated dataset
         ($N = 1{,}954$).}
\label{tab:gene_dist}
\begin{tabular}{@{}lrr@{}}
\toprule
\textbf{Gene} & \textbf{Variants ($n$)} & \textbf{Fraction (\%)} \\
\midrule
\textit{MYH7}    & 986 & 50.5 \\
\textit{MYBPC3}  & 429 & 22.0 \\
\textit{TNNT2}   & 130 &  6.7 \\
\textit{TNNI3}   & 118 &  6.0 \\
\textit{TPM1}    & 100 &  5.1 \\
\textit{ACTC1}   &  64 &  3.3 \\
\textit{MYL2}    &  54 &  2.8 \\
\textit{MYL3}    &  37 &  1.9 \\
\textit{TNNC1}   &  36 &  1.8 \\
\midrule
\textbf{Total}   & \textbf{1,954} & \textbf{100.0} \\
\bottomrule
\end{tabular}
\end{table}

The resulting class distribution is \textbf{73.7\,\% Pathogenic}
(1,441~variants) versus \textbf{26.3\,\% Benign} (513~variants). This
imbalance reflects genuine ascertainment bias in clinical databases—disease
cohorts are sequenced far more often than control populations—and must be
handled explicitly in model training.

\subsection{Feature Engineering: The 51-Feature Schema}
\label{subsec:features}

Fifty-one features were engineered to represent each variant from four
conceptually distinct perspectives: (1)~the \emph{physicochemical}
consequence of the amino-acid substitution itself; (2)~the \emph{structural
and functional context} of the mutation site within the protein; (3)~the
\emph{evolutionary constraint} at that site; and (4)~a \emph{gene-identity}
prior. Table~\ref{tab:features} summarises the four groups.

\begin{table}[H]
\centering
\caption{The 51-feature schema grouped by biological rationale.
         Features fed to the XGBoost tabular stream: 29.
         Sequence window features fed to the CNN stream: 22.}
\label{tab:features}
\begin{tabular}{@{}lp{6.5cm}r@{}}
\toprule
\textbf{Category} & \textbf{Features} & \textbf{Count} \\
\midrule
\multirow{4}{*}{\shortstack[l]{\textbf{Physicochemical}\\(substitution-level)}}
  & Grantham score (composite of composition, polarity, molecular volume) & \\
  & Reference and alternate residue size, charge & \\
  & Size change $\Delta s$, charge change $\Delta q$ & \\
  & & 7 \\
\midrule
\multirow{5}{*}{\shortstack[l]{\textbf{Structural /}\\{\textbf{Functional}}}}
  & Binary indicators: functional domain, coiled-coil, $\alpha$-helix,
    $\beta$-strand, turn, disordered region, compositionally biased region,
    post-translational modification (PTM) site, functional site,
    secondary structure presence, named region & \\
  & & 11 \\
\midrule
\multirow{2}{*}{\shortstack[l]{\textbf{Evolutionary /}\\{\textbf{Positional}}}}
  & Relative position within protein (normalised $[0,1]$) & \\
  & Population allele frequency (gnomAD/ExAC/TOPMed) & 2 \\
\midrule
\textbf{Gene Identity}
  & One-hot encoding: \textit{MYH7}, \textit{MYBPC3}, \textit{TNNT2},
    \textit{TNNI3}, \textit{TPM1}, \textit{MYL2}, \textit{MYL3},
    \textit{ACTC1}, \textit{TNNC1} & 9 \\
\midrule
\multirow{2}{*}{\shortstack[l]{\textbf{Sequence Window}\\(CNN only)}}
  & Residue size category and charge category at each of the 11 positions
    in the $[\text{pos}-5,\, \text{pos}+5]$ sliding window & 22 \\
\midrule
\multicolumn{2}{@{}l}{\textbf{Total}} & \textbf{51} \\
\bottomrule
\end{tabular}
\end{table}

\paragraph{Why the Grantham Score?}
The Grantham score~\cite{grantham1974} quantifies the expected impact of an
amino-acid substitution along three orthogonal biochemical axes: atomic
composition, polarity, and molecular volume. A glycine-to-arginine
substitution (Grantham $= 125$) introduces a large, positively charged side
chain into a position previously occupied by the simplest amino acid,
creating a severe steric and electrostatic mismatch. Empirically, the
pathogenic variants in this dataset have a mean Grantham score of
\textbf{79.78} compared to \textbf{60.85} for Benign variants
(Table~\ref{tab:grantham}), confirming that radical substitutions are
systematically enriched in pathogenic calls.

\begin{table}[H]
\centering
\caption{Mean Grantham score by pathogenicity class.}
\label{tab:grantham}
\begin{tabular}{@{}lcc@{}}
\toprule
\textbf{Class} & \textbf{$n$} & \textbf{Mean Grantham Score} \\
\midrule
Pathogenic & 1,441 & 79.78 \\
Benign     &   513 & 60.85 \\
\bottomrule
\end{tabular}
\end{table}

\paragraph{Structural Annotation Features.}
UniProt provides rich, expert-curated positional annotations for all nine
HCM proteins. Binary indicators encode whether a variant falls within a
known functional domain (e.g., the myosin motor domain of MYH7), a
coiled-coil region, an alpha-helix, a beta-strand, a disordered linker, a
PTM site, or a characterised functional site. These indicators encode
knowledge that is invisible to sequence-only methods but is directly
relevant to structural stability: a mutation at an ATP-binding site or a
known phosphorylation residue carries an immediate functional implication,
regardless of its Grantham score.

\paragraph{Positional and Population Features.}
The relative position of a variant within its protein (normalised to
$[0,1]$) allows the model to learn, without explicit structural coordinates,
that the myosin head domain of MYH7 is far more intolerant of substitution
than the distal tail. The population allele frequency from gnomAD/ExAC/TOPMed
provides evolutionary evidence: variants observed at frequencies above
0.1\,\% in healthy control populations are extremely unlikely to be
fully penetrant Mendelian disease alleles.

\subsection{Sliding-Window Sequence Extraction}
\label{subsec:window}

The 11-residue sliding window centred on the mutated position
(positions $-5$ to $+5$) was extracted from the full canonical protein
sequence retrieved from UniProt for each accession. At each of the 11
positions, two categorical encodings were stored: a \emph{size category}
(small/medium/large mapped to integer values 1/2/3) and a \emph{charge
category} (negative/$-1$, neutral/$0$, positive/$+1$). This yields a
two-channel sequence tensor of shape $(11, 2)$ per variant.

Boundary variants at the protein N- or C-terminus where the full
five-residue flanking context does not exist are zero-padded
(size $= 0$, charge $= 0$), consistent with a neutral encoding that
introduces no spurious physicochemical signal. The window is centred on the
\emph{reference} residue; the convolutional branch therefore learns what
physico\-chemical neighbourhood is \emph{expected} at that position, and
the auxiliary features communicate the \emph{disruption} introduced by the
alternate residue.

\begin{figure}[H]
\centering
% Replace with actual sliding window schematic
\fbox{\parbox{0.8\textwidth}{\centering\vspace{2cm}
\textbf{[FIGURE PLACEHOLDER]}\\
Sliding-window diagram: 11-residue window centred on mutation site,
showing size (S) and charge (Q) encoding for each position.
Zero-padding at protein boundaries is indicated in grey.
\vspace{2cm}}}
\caption{Schematic of the 11-residue sliding window sequence extraction.
The mutated position is highlighted at the centre; size and charge
encodings are shown for each flanking position.}
\label{fig:window}
\end{figure}

\subsection{The Three-Branch Comparative Experiment}
\label{subsec:three_branch}

The core experimental design of Phase~II is a controlled three-way
comparison, structured to isolate the individual contributions of tabular
expert knowledge and sequence-derived context, and then to quantify the
benefit of fusing them.

\begin{description}[leftmargin=2em, labelwidth=1.8em]

\item[\textbf{Branch 1 — Tuned XGBoost (Tabular):}]
A gradient-boosted decision tree ensemble trained exclusively on the
29-feature tabular vector (physicochemical, structural, evolutionary, and
gene-identity features). Hyperparameters were tuned via grid search. This
branch represents the \emph{expert knowledge baseline}—the maximum
performance achievable from hand-crafted, interpretable features alone.

\item[\textbf{Branch 2 — 1D-CNN (Sequence + Auxiliary):}]
A multi-input neural network processing the $(11 \times 2)$ sequence tensor
through a convolutional branch in parallel with the 29-feature auxiliary
vector through a dense MLP branch, before fusing and outputting a sigmoid
pathogenicity probability. This branch tests whether local sequence context
provides discriminative signal above the tabular baseline.

\item[\textbf{Branch 3 — Late-Fusion Hybrid Ensemble:}]
A Logistic Regression meta-learner trained on out-of-fold (OOF) probability
estimates from Branches~1 and~2. This stacking layer learns the optimal
linear combination of the two models' confidence scores. Any improvement in
Branch~3 over Branches~1 and~2 can be attributed specifically to information
complementarity between modalities.

\end{description}

\subsection{1D-CNN Architecture}
\label{subsec:cnn}

The 1D-CNN was implemented in TensorFlow/Keras. Table~\ref{tab:cnn} details
the exact layer configuration extracted from the trained model artefact.

\begin{table}[H]
\centering
\caption{Exact architecture of the 1D-CNN used in Branch 2. The model
         has two input branches fused before the final classification head.}
\label{tab:cnn}
\begin{tabular}{@{}llll@{}}
\toprule
\textbf{Branch} & \textbf{Layer} & \textbf{Configuration} & \textbf{Output Shape} \\
\midrule
\multirow{3}{*}{Sequence}
  & Conv1D             & 32 filters, kernel=3, padding=`same', ReLU & $(11, 32)$ \\
  & BatchNormalization & —                                           & $(11, 32)$ \\
  & GlobalMaxPooling1D & —                                           & $(32,)$ \\
\midrule
\multirow{4}{*}{Auxiliary}
  & Dense              & 64 units, ReLU                              & $(64,)$ \\
  & Dropout            & rate $= 0.3$                                & $(64,)$ \\
  & Dense              & 32 units, ReLU                              & $(32,)$ \\
  & Dropout            & rate $= 0.3$                                & $(32,)$ \\
\midrule
\multirow{4}{*}{Fusion Head}
  & Concatenate        & CNN output $(32)$ + MLP output $(32)$       & $(64,)$ \\
  & Dense              & 32 units, ReLU                              & $(32,)$ \\
  & Dropout            & rate $= 0.3$                                & $(32,)$ \\
  & Dense              & 16 units, ReLU                              & $(16,)$ \\
  & Dense              & 1 unit, Sigmoid                             & $(1,)$ \\
\bottomrule
\end{tabular}
\end{table}

The choice of kernel size 3 in the convolutional layer is deliberate: a
window of three consecutive residue encodings captures tripeptide-level
physicochemical motifs—the smallest structural unit that encodes directional
secondary-structure preference. Batch Normalisation after the convolutional
layer stabilises the activation distribution across folds and prevents
internal covariate shift in a dataset of this modest size.
GlobalMaxPooling1D then extracts the single strongest activation across all
11 window positions, providing a compact, position-invariant summary of the
most informative local motif.

\subsection{Training Protocol and Class Imbalance Handling}
\label{subsec:training}

\paragraph{5-Fold Stratified Cross-Validation.}
All model evaluation was performed under \textbf{5-Fold Stratified
Cross-Validation} (StratifiedKFold, scikit-learn). Stratification ensures
that each of the five validation folds maintains the global 73.7\,\%/26.3\,\%
class ratio, preventing evaluation folds from being dominated by the majority
class. Without stratification, a single unlucky fold containing only 15--20
Benign variants would yield inflated accuracy metrics without genuinely
testing the model's ability to recognise Benign variants.

The meta-learner (Branch~3) was additionally trained on OOF probability
estimates accumulated over the five folds, a canonical stacking protocol
that prevents data leakage from the base models into the meta-model's
training set.

\paragraph{Class Imbalance Mitigation.}
The Logistic Regression meta-learner used
\texttt{class\_weight=`balanced'}, which applies inverse-frequency sample
weights $(w_i = N / (2 \cdot n_c)$ for class $c)$, penalising
misclassification of the minority Benign class more heavily. The XGBoost
base model used a tuned \texttt{scale\_pos\_weight} parameter. The CNN
training used a weighted binary cross-entropy loss. The operating
classification threshold was tuned by maximising Youden's J statistic on
OOF probabilities, yielding an optimal threshold of \textbf{0.5294}.

%=======================================================================
\section{Experimental Results and Comparative Analysis}
\label{sec:results}

\subsection{Performance Audit: Three-Model Comparison}
\label{subsec:performance}

Table~\ref{tab:perf} presents the head-to-head performance comparison of
all three experimental branches at the optimised classification threshold,
evaluated under 5-Fold Stratified Cross-Validation.

\begin{table}[H]
\centering
\caption{Comparative performance of the three experimental branches at
         the optimal Youden threshold (0.5294 for Ensemble).
         Bold values indicate the best result per column.}
\label{tab:perf}
\begin{tabular}{@{}lcccccc@{}}
\toprule
\textbf{Model} & \textbf{ROC-AUC} & \textbf{PR-AUC} & \textbf{Accuracy}
  & \textbf{Bal.\ Acc.} & \textbf{F1-Macro} & \textbf{MCC} \\
\midrule
Tuned XGBoost
  & 0.7402 & 0.8639 & \textbf{0.7702} & 0.5818 & 0.5800 & 0.293 \\
1D-CNN
  & 0.7934 & 0.9049 & 0.7288 & 0.7257 & 0.6906 & 0.408 \\
\textbf{Hybrid Ensemble}
  & \textbf{0.8222} & \textbf{0.9141} & 0.7421
  & \textbf{0.7529} & \textbf{0.7095} & \textbf{$\approx$0.45} \\
\bottomrule
\end{tabular}
\end{table}

\noindent
The Matthews Correlation Coefficient is defined as:
\begin{equation}
MCC = \frac{TP \cdot TN - FP \cdot FN}
           {\sqrt{(TP+FP)(TP+FN)(TN+FP)(TN+FN)}}
\label{eq:mcc}
\end{equation}
where $TP$, $TN$, $FP$, $FN$ are the true positive, true negative, false
positive, and false negative counts, respectively. MCC is preferred over
raw accuracy as the primary metric here because it is invariant to class
imbalance and penalises failure on both classes symmetrically.

\subsection{Hybrid Model Confusion Matrix}
\label{subsec:cm}

Table~\ref{tab:cm} presents the confusion matrix of the Hybrid Ensemble at
the optimal threshold of 0.5294.

\begin{table}[H]
\centering
\caption{Confusion matrix of the Late-Fusion Hybrid Ensemble
         ($N = 1{,}954$, optimal threshold $= 0.5294$).}
\label{tab:cm}
\begin{tabular}{@{}llcc@{}}
\toprule
& & \multicolumn{2}{c}{\textbf{Predicted}} \\
\cmidrule(l){3-4}
& & \textbf{Benign (0)} & \textbf{Pathogenic (1)} \\
\midrule
\multirow{2}{*}{\textbf{Actual}}
  & \textbf{Benign (0)}     & 398 (TN) & 115 (FP) \\
  & \textbf{Pathogenic (1)} & 389 (FN) & 1052 (TP) \\
\bottomrule
\end{tabular}
\end{table}

\noindent
From the confusion matrix:
\noindent
Treating each class in turn as the ``positive'' class:
\begin{align}
\text{Recall (Benign)}
  &= \frac{\text{correctly predicted Benign}}{\text{all actual Benign}}
   = \frac{TN}{TN + FP} = \frac{398}{398+115} = 0.776 \label{eq:rec_b}\\
\text{Recall (Pathogenic)}
  &= \frac{\text{correctly predicted Pathogenic}}{\text{all actual Pathogenic}}
   = \frac{TP}{TP + FN} = \frac{1052}{1052+389} = 0.730 \label{eq:rec_p}\\
\text{Precision (Benign)}
  &= \frac{\text{correctly predicted Benign}}{\text{all predicted Benign}}
   = \frac{TN}{TN + FN} = \frac{398}{398+389} = 0.506 \label{eq:prec_b}
\end{align}

\subsection{The Hybrid Advantage}
\label{subsec:advantage}

The most revealing contrast in Table~\ref{tab:perf} is between the
XGBoost model and the Hybrid. The XGBoost achieves the \emph{highest raw
accuracy} (77.0\,\%) yet the \emph{lowest MCC} (0.293) and the lowest
Balanced Accuracy (0.582)—a pattern that unmistakably identifies
majority-class collapse. Inspection of the class-specific recalls confirms
this: XGBoost achieves Benign recall of only \textbf{18.5\,\%} while
correctly classifying 97.8\,\% of Pathogenic variants. In practical terms,
this model would incorrectly flag roughly four out of every five genuinely
Benign variants as disease-causing—an unacceptable clinical performance.

The 1D-CNN corrects this dramatically, raising Benign recall to
\textbf{71.9\,\%} at the cost of a modest reduction in Pathogenic recall to
73.2\,\%. The CNN achieves this by encoding local sequence context: it
learns that certain physico\-chemical neighbourhood patterns are
characteristic of tolerated polymorphisms, even when the tabular features
appear moderately suspicious. The MCC improvement of $+0.115$ over XGBoost
reflects this genuine gain in discriminative balance.

The Hybrid Ensemble then delivers a further $+0.042$ MCC improvement by
learning that the two probability streams carry \emph{non-redundant
information}. When the XGBoost is highly confident about a Pathogenic call
(high Grantham score, disrupted functional domain) but the CNN is
uncertain, and vice versa, the meta-learner's logistic decision boundary
correctly combines the two signals. The net result is a 4-percentage-point
improvement in Benign recall over the CNN ($77.6\,\%$ vs.\
$71.9\,\%$) while maintaining comparable Pathogenic recall ($73.0\,\%$
vs.\ $73.2\,\%$).

\begin{figure}[H]
\centering
\begin{subfigure}[b]{0.48\textwidth}
  \fbox{\parbox{\textwidth}{\centering\vspace{2cm}
  \textbf{[FIGURE PLACEHOLDER]}\\
  ROC curve comparison: XGBoost (AUC=0.74), CNN (AUC=0.79),
  Hybrid (AUC=0.82). Insert \texttt{roc\_comparison.png} here.
  \vspace{2cm}}}
  \caption{ROC curves for all three experimental branches.}
  \label{fig:roc}
\end{subfigure}
\hfill
\begin{subfigure}[b]{0.48\textwidth}
  \fbox{\parbox{\textwidth}{\centering\vspace{2cm}
  \textbf{[FIGURE PLACEHOLDER]}\\
  Precision-Recall curve comparison. Insert \texttt{pr\_comparison.png}
  here.
  \vspace{2cm}}}
  \caption{Precision-Recall curves. PR-AUC is more sensitive to class
           imbalance than ROC-AUC.}
  \label{fig:pr}
\end{subfigure}
\caption{Performance curves comparing the three model branches under
         5-Fold Stratified Cross-Validation. The Hybrid Ensemble
         dominates in both ROC and PR space.}
\label{fig:curves}
\end{figure}

\begin{figure}[H]
\centering
\fbox{\parbox{0.6\textwidth}{\centering\vspace{2.5cm}
\textbf{[FIGURE PLACEHOLDER]}\\
Confusion matrix heatmap for the Hybrid Ensemble.
Insert \texttt{confusion\_matrix.png} here.
\vspace{2.5cm}}}
\caption{Confusion matrix of the Late-Fusion Hybrid Ensemble at the
         optimal threshold (0.5294). Colour intensity encodes cell count.}
\label{fig:cm}
\end{figure}

\subsection{Error Analysis and Dataset Limitations}
\label{subsec:errors}

The 389 false-negative Pathogenic variants (missed by the Hybrid) represent
the model's primary failure mode. These are concentrated among variants with
moderate Grantham scores ($50$--$80$), low population frequency, and sparse
structural annotation — the ``intermediate'' cases where human expert
annotation is itself uncertain. Two structural factors amplify this
limitation:

\begin{enumerate}
\item \textbf{Gene under-representation.} Four genes—\textit{ACTC1} (n=64),
  \textit{MYL2} (n=54), \textit{MYL3} (n=37), and \textit{TNNC1} (n=36)—
  together contribute fewer than 200 variants to the training set.
  For these genes, UniProt structural annotations are sparser, population
  frequency data is limited due to fewer published cohort studies, and the
  model's gene-specific priors are learned from inadequate statistical
  support.

\item \textbf{Sequence window sparsity at protein termini.} Variants
  within the first or last five residues of a protein receive zero-padded
  flanking context, reducing the discriminative information available to
  the CNN branch. For these edge-position variants, the model relies almost
  entirely on the tabular branch.
\end{enumerate}

The 115 false-positive Benign variants (incorrectly flagged as Pathogenic)
are particularly costly in a clinical context. These variants carry
relatively high Grantham scores or fall within annotated functional domains,
leading the tabular model to assign elevated pathogenicity probabilities that
the CNN cannot fully offset.

%=======================================================================
\section{Interpretability and Biological Logic}
\label{sec:interp}

\subsection{XGBoost SHAP Feature Importance Hierarchy}
\label{subsec:shap}

SHAP (SHapley Additive exPlanations) analysis of the XGBoost model
identifies a biologically coherent hierarchy of predictive features,
summarised in Table~\ref{tab:shap}.

\begin{table}[H]
\centering
\caption{Top-5 features by mean absolute SHAP value in the tuned XGBoost
         model, with biological interpretation.}
\label{tab:shap}
\begin{tabular}{@{}clp{7cm}@{}}
\toprule
\textbf{Rank} & \textbf{Feature} & \textbf{Biological Rationale} \\
\midrule
1 & Grantham Score
  & Composite measure of substitution radicality. Higher scores indicate
    physicochemical incompatibility between reference and alternate residues,
    predicting structural disruption. \\
2 & Relative Position
  & Mutations in the N-terminal motor domain of MYH7 ($\text{position} < 0.4$)
    show systematically elevated pathogenicity; the model learns the protein's
    structural topology implicitly. \\
3 & Functional/PTM Site Indicator
  & Variants disrupting ATP-binding sites, actin-binding interfaces, or
    regulatory phosphorylation residues have near-certain pathogenicity,
    regardless of Grantham score. \\
4 & Domain Indicator
  & Presence within a characterised structural domain (e.g., myosin-binding
    domain of MYBPC3, EF-hand domain of TNNC1) elevates pathogenicity
    probability due to reduced substitution tolerance in folded domain
    interiors. \\
5 & Population Frequency
  & Strong \emph{negative} predictor: variants observed at allele frequency
    $> 0.001$ in gnomAD are highly unlikely to be fully penetrant HCM alleles;
    the model applies a large protective Benign weight to such variants. \\
\bottomrule
\end{tabular}
\end{table}

\begin{figure}[H]
\centering
\fbox{\parbox{0.8\textwidth}{\centering\vspace{2.5cm}
\textbf{[FIGURE PLACEHOLDER]}\\
SHAP summary beeswarm plot: feature-level contribution to XGBoost
output across all 1,954 variants. Generate using
\texttt{shap.summary\_plot(shap\_values, X\_tabular)}.
\vspace{2.5cm}}}
\caption{SHAP summary plot for the Tuned XGBoost model. Each point
         represents one variant; colour encodes feature value (red=high,
         blue=low). The Grantham score and relative position dominate
         positive SHAP contributions for Pathogenic calls.}
\label{fig:shap}
\end{figure}

\subsection{CNN Sequence-Context Insights}
\label{subsec:cnn_interp}

The 1D-CNN's substantial improvement in Benign recall
($71.9\,\%$ vs.\ $18.5\,\%$ for XGBoost) provides indirect evidence about
the discriminative content of the 11-residue window. While explicit gradient
saliency maps were not generated in the current training run, several
inferences can be drawn from the model's performance profile.

The convolutional filters (kernel size 3) operate across triplets of
consecutive $(size, charge)$ encodings. Filters that activate maximally on
patterns such as $(large{+}, large{+}, large{+})$ or $(neutral, neutral,
negative{-})$ learn to detect regions of consistent physicochemical
character — the fingerprint of structured, evolutionarily constrained
secondary structures. A mutation inserted into such a context with a
radically different physicochemical profile (as captured by the Grantham
score and delta features in the auxiliary branch) creates a detectable
\emph{neighbourhood mismatch}.

The GlobalMaxPooling layer retains only the strongest such activation across
all 11 positions, suggesting that the model's sequence-level decision is
driven by the \emph{most anomalous tripeptide motif} in the window rather
than by an average neighbourhood property. This aligns with the biological
principle that protein instability from point mutations is often localised
to a single critical interaction rather than distributed evenly across the
neighbourhood.

\begin{figure}[H]
\centering
\fbox{\parbox{0.8\textwidth}{\centering\vspace{2.5cm}
\textbf{[FIGURE PLACEHOLDER]}\\
Saliency map / integrated gradients heatmap: per-position importance
in the 11-residue window, averaged over all correctly-classified
Pathogenic variants. Positions $-1$ and $+1$ (immediately flanking the
mutation site) are expected to show the highest saliency.
\vspace{2.5cm}}}
\caption{Gradient-based saliency map for the CNN sequence branch,
         showing per-position contribution to the pathogenicity prediction
         averaged over all correctly classified Pathogenic variants.}
\label{fig:saliency}
\end{figure}

\subsection{Biological Conclusion: Model Mathematics Meets Cardiac Reality}
\label{subsec:bio}

The model's decision hierarchy—Grantham score, relative position, functional
site disruption, domain membership, population frequency—mirrors the clinical
geneticist's mental model for evaluating HCM variants. What the machine
learning framework adds is a \emph{quantified, consistent, and scalable}
instantiation of that model, immune to cognitive fatigue and applicable at
genomic scale.

The CNN's contribution is particularly compelling in biological terms. The
sarcomeric proteins are not random polypeptides: they form precisely
organised, mechanically stressed macromolecular assemblies in which every
residue's identity within its local neighbourhood is under strong purifying
selection. A point mutation that locally disrupts the charge pattern of an
alpha-helix in the myosin S1 domain does not merely affect a single
hydrogen bond; it propagates mechanical instability through the entire
force-transmission apparatus, impairing crossbridge cycling and ultimately
reducing cardiac contractile efficiency. The CNN's neighbourhood encoding
captures this collective structural sensitivity in a mathematically compact
representation that is inaccessible to feature-by-feature tabular analysis.

The Grantham mean difference ($79.78$ versus $60.85$, a gap of
approximately \textbf{19 Grantham units}) confirms that pathogenic HCM
mutations involve systematically more radical amino-acid substitutions than
Benign polymorphisms. This single observation validates the biological
logic underlying the entire feature engineering strategy.

%=======================================================================
\section{Discussion and Future Roadmap}
\label{sec:discussion}

\subsection{Phase II Limitations}
\label{subsec:limits}

Despite the Hybrid model's strong performance, three categories of limitation
constrain the generalisability of Phase~II findings.

\paragraph{Gene-Specific Data Sparsity.}
The four smallest gene cohorts—\textit{ACTC1}, \textit{MYL2},
\textit{MYL3}, and \textit{TNNC1}—are substantially under-powered for
supervised learning. With fewer than 65 labelled variants per gene, the
model's gene-specific priors (encoded by one-hot gene identity features)
are learned from very limited statistical support. Gene-specific subgroup
analysis would likely reveal substantially lower performance for these genes
relative to \textit{MYH7} and \textit{MYBPC3}. Acquisition of additional
labelled variants for rare-gene HCM subtypes is a priority for Phase~III
data enrichment.

\paragraph{VUS Exclusion.}
The deliberate exclusion of VUS variants from training, while scientifically
rigorous for supervised learning, means that the model's output for novel
VUS cases during inference is extrapolated beyond the training distribution.
Phase~III should include a held-out VUS validation set with expert consensus
labels to quantify extrapolation confidence.

\paragraph{Absence of Protein 3D Structural Coordinates.}
The current feature set relies on sequence-derived physicochemical proxies
rather than three-dimensional structural data. Variants that are spatially
proximal in the folded protein but distant in primary sequence — a common
feature of allosteric HCM mutations — cannot be captured by the current
encoding. Integration of AlphaFold2 structural coordinates is planned for
Phase~III.

\paragraph{Negative Predictive Value.}
The Hybrid model's Precision for the Benign class is \textbf{0.506}
(Eq.~\ref{eq:prec_b}), meaning that approximately half of variants
predicted as Benign are actually Pathogenic. This limits the model's utility
as a standalone ``rule-out'' tool in clinical settings without additional
orthogonal evidence.

\subsection{Phase III Roadmap: External Validation and Architecture Expansion}
\label{subsec:phase3}

Phase~III will extend the Phase~II framework along two parallel tracks:

\begin{description}[leftmargin=2em]

\item[\textbf{Track A — External Validation on SHaRe Registry:}]
The Sarcomeric Human Cardiomyopathy Registry (SHaRe) provides a large,
independently collected, multi-centre cohort of HCM probands with
detailed genotypic and phenotypic data. Variants in the SHaRe dataset
that were not present in ClinVar at the time of Phase~II dataset
construction will be used as a fully held-out external validation set.
This will test whether the model generalises beyond its training
distribution and whether its probability calibration (Figure placeholder
for \texttt{calibration\_plot.png}) is reliable in a prospective clinical
setting.

\item[\textbf{Track B — ESM-2 Protein Language Model Embeddings:}]
The 11-residue $(size, charge)$ encoding used in Phase~II, while
interpretable, is a coarse discretisation of the amino-acid alphabet.
Meta AI's ESM-2 protein language model provides 480-dimensional or
1,280-dimensional per-residue embeddings pre-trained on 250 million
protein sequences, capturing evolutionary co-variation, predicted secondary
structure, and long-range residue contacts implicitly within a single
vector. Replacing the manual window encoding with ESM-2 embeddings
is expected to substantially enrich the sequence information available to
the CNN branch and improve discrimination in the under-annotated
small-gene cohorts. Specifically, for \textit{ACTC1} and \textit{MYL2}
variants where UniProt structural annotations are sparse, ESM-2 embeddings
encode cross-family evolutionary knowledge that may compensate for the
limited within-gene training data.

\end{description}

%=======================================================================
\section{Conclusion}
\label{sec:conclusion}

Phase~II of this RVCE Experiential Learning project has delivered a
rigorously validated comparative analysis of tabular and sequence-based
machine learning approaches for HCM variant pathogenicity classification.
Three key findings emerge:

\begin{enumerate}
\item \textbf{Tabular features alone are insufficient.} The tuned XGBoost
  model, despite encoding rich biochemical and structural knowledge,
  collapses to majority-class prediction under the 73.7\,\%/26.3\,\% class
  imbalance, achieving a Benign recall of only 18.5\,\% (MCC $= 0.293$).

\item \textbf{Sequence context restores minority-class discrimination.} The
  1D-CNN's local window encoding raises Benign recall to 71.9\,\% and MCC
  to $0.408$—a 39\,\% relative improvement—by capturing neighbourhood-level
  physicochemical patterns that distinguish tolerated polymorphisms from
  pathogenic mutations.

\item \textbf{Hybrid fusion is additive.} The Late-Fusion Ensemble achieves
  the best performance on all balanced metrics: ROC-AUC $= 0.822$,
  Balanced Accuracy $= 0.753$, F1-Macro $= 0.710$, and MCC $\approx 0.45$,
  confirming the complementarity hypothesis.

\end{enumerate}

These results validate the core design philosophy of Phase~II: that expert
biochemical knowledge and learned sequence-context patterns are genuinely
complementary information sources, and that their fusion via a late-stacking
meta-learner extracts measurable additional discriminative value.

The primary limitation—low Benign precision ($0.506$) and the 389
false-negative pathogenic variants—defines the central challenge for
Phase~III: improving sensitivity for the difficult intermediate-confidence
variants through richer sequence representations (ESM-2) and expanded
training data from the SHaRe registry. The transition from Phase~II
to Phase~III represents a shift from \emph{proof-of-concept hybrid
validation} to \emph{clinical-grade generalisation testing}, a necessary
step on the path toward a deployable decision-support tool for HCM genetic
counselling.

%=======================================================================
% References
%=======================================================================
\begin{thebibliography}{99}

\bibitem{grantham1974}
Grantham, R. (1974).
Amino acid difference formula to help explain protein evolution.
\textit{Science}, 185(4154), 862--864.
\texttt{doi:10.1126/science.185.4154.862}

\bibitem{clinvar}
Landrum, M. J., et al. (2016).
ClinVar: public archive of interpretations of clinically relevant variants.
\textit{Nucleic Acids Research}, 44(D1), D862--D868.

\bibitem{gnomad}
Karczewski, K. J., et al. (2020).
The mutational constraint spectrum quantified from variation in 141,456 humans.
\textit{Nature}, 581, 434--443.

\bibitem{acmg}
Richards, S., et al. (2015).
Standards and guidelines for the interpretation of sequence variants: a joint
consensus recommendation of the American College of Medical Genetics and Genomics
and the Association for Molecular Pathology.
\textit{Genetics in Medicine}, 17(5), 405--424.

\bibitem{esm2}
Lin, Z., et al. (2023).
Evolutionary-scale prediction of atomic-level protein structure with a language model.
\textit{Science}, 379(6637), 1123--1130.

\bibitem{alphafold}
Jumper, J., et al. (2021).
Highly accurate protein structure prediction with AlphaFold.
\textit{Nature}, 596, 583--589.

\bibitem{uniprot}
The UniProt Consortium. (2021).
UniProt: the universal protein knowledgebase in 2021.
\textit{Nucleic Acids Research}, 49(D1), D480--D489.

\bibitem{xgboost}
Chen, T., \& Guestrin, C. (2016).
XGBoost: A scalable tree boosting system.
\textit{Proceedings of the 22nd ACM SIGKDD International Conference on
Knowledge Discovery and Data Mining}, 785--794.

\bibitem{share}
Ho, C. Y., et al. (2018).
Genotype and lifetime burden of disease in hypertrophic cardiomyopathy:
insights from the Sarcomeric Human Cardiomyopathy Registry (SHaRe).
\textit{Circulation}, 138(14), 1387--1398.

\bibitem{shap}
Lundberg, S. M., \& Lee, S.-I. (2017).
A unified approach to interpreting model predictions.
\textit{Advances in Neural Information Processing Systems}, 30.

\end{thebibliography}

\end{document}
%=======================================================================
% END OF DOCUMENT
%=======================================================================
